% ============================================================================
% I. Видеофайлът като обект на експертиза
% ============================================================================
\section{Видеофайлът като обект на експертиза}

Един видеофайл съдържа две различни нива на структура. Външното ниво е
\term{контейнерът}: той описва какви потоци има във файла (видео, звук, субтитри),
къде във файла се намира всяка порция данни и кога трябва да бъде възпроизведена.
Вътрешното ниво е \term{кодекът}: той определя как самото изображение е компресирано
в поредица от байтове. За възстановяването двете нива имат различно значение. Загубата
на контейнера прави файла нечетим за обикновен плейър, но компресираните кадри може
да са още налице. Загубата на част от кодираните данни поврежда самите кадри.

В работата е разгледана най-разпространената комбинация: контейнер MP4 и видеокодек
H.264 (наричан още AVC). Изборът не е случаен. MP4 и сродният му формат MOV се
използват масово от мобилни телефони и камери, а H.264 е кодекът, с който работи
голяма част от записващите устройства.

\subsection{Контейнерът MP4 и базовият медиен формат на ISO}

Форматът MP4 \cite{iso14496_14} е частен случай на \term{базовия медиен файлов формат
на ISO} (ISO Base Media File Format, ISOBMFF), дефиниран в ISO/IEC 14496-12
\cite{iso14496_12}. Файлът е изграден от последователност от \term{кутии} (boxes, в
по-старите текстове \term{atoms}). Всяка кутия започва със заглавие от два елемента:
размер от 4 байта и тип от 4 знака, например \code{ftyp}, \code{moov} или \code{mdat}.
Ако размерът е 1, след типа следва 64-битов размер. Ако размерът е 0, кутията
продължава до края на файла. Кутиите могат да се влагат една в друга, така че файлът
образува дърво.

Три кутии от най-горно ниво са важни за възстановяването \figref{fig:mp4}:

\begin{itemize}[leftmargin=2cm]
  \item \code{ftyp} (file type) посочва към кой вариант на формата принадлежи файлът
        и се поставя възможно най-рано в него. Когато е първата кутия, типът
        \code{ftyp} стои на байтове 4 до 7 от началото на файла и служи като
        \term{сигнатура} при търсене на изтрити файлове.
  \item \code{mdat} (media data) съдържа самите компресирани данни на видеото и
        звука. Това обикновено е най-голямата част от файла.
  \item \code{moov} (movie) съдържа метаданните: описание на всяка писта
        (\code{trak}) и таблиците на пробите (\code{stbl}).
\end{itemize}

\begin{figure}[H]
\centering
\begin{tikzpicture}[node distance=0.35cm]
  \node[block, minimum width=3.0cm] (ftyp) {\code{ftyp}\\ тип на файла};
  \node[accent, right=of ftyp, minimum width=4.6cm] (moov) {\code{moov}\\ метаданни};
  \node[danger, right=of moov, minimum width=4.6cm] (mdat) {\code{mdat}\\ компресирани данни};

  \node[note, below=0.9cm of moov.south east, anchor=north east, text width=6.6cm] (stbl)
    {\code{trak} $\rightarrow$ \code{mdia} $\rightarrow$ \code{minf} $\rightarrow$ \code{stbl}:\\
     \code{stsd} (параметри на кодека, SPS и PPS)\\
     \code{stsz} (размер на всяка проба)\\
     \code{stco} / \code{co64} (отместване на всяко парче)\\
     \code{stss} (кои проби са ключови кадри)};
  \node[note, right=1.4cm of stbl, text width=4.0cm] (nal)
    {Пробите на видеото:\\
     {}[дължина][NAL единица]\\
     {}[дължина][NAL единица]\\
     \ldots};

  \draw[line] (moov.south) to (moov.south |- stbl.north);
  \draw[line] (mdat.south) to (mdat.south |- nal.north);
  \draw[line, dashed] (stbl.east) to node[above, font=\scriptsize]{адреси} (nal.west);
\end{tikzpicture}
\caption{Опростена структура на файл MP4 (по \cite{iso14496_12, iso14496_15})}
\label{fig:mp4}
\end{figure}

Ключовото наблюдение е, че кутията \code{mdat} сама по себе си не е самоописваща се.
Тя е непрекъснат поток от байтове. Къде започва и свършва всеки кадър, личи само от
таблиците в \code{moov}: \code{stsz} дава размера на всяка проба, \code{stco} или
\code{co64} дава отместването на всяко парче (chunk) от началото на файла, а
\code{stsc} свързва пробите с парчетата \cite{iso14496_12}. Параметрите, нужни за
декодиране на видеото, също са в \code{moov}, в описанието на пробите \code{stsd}.
Ако \code{moov} липсва, плейърът не знае нито къде са кадрите, нито как да ги декодира.

Стандартът предвижда и \term{фрагментиран} вариант, при който метаданните са
разпределени в множество кутии \code{moof}, всяка от които описва следващия я участък
от данни \cite{iso14496_12}. Така файлът остава четим до последния завършен фрагмент.
Значението на тази разлика за повредени записи е разгледано в раздел~II.

\subsection{Видеопотокът H.264}

Стандартът H.264 \cite{itu_h264} дели кодирания поток на \term{единици на слоя за
мрежова абстракция} (Network Abstraction Layer units, NAL). Всяка NAL единица започва
с 1 байт заглавие. В него 5 бита задават типа на единицата, 2 бита (\code{nal\_ref\_idc})
показват дали тя служи за опора при декодиране на други кадри, а 1 бит трябва винаги
да е нула. Таблица~\ref{tab:nal} изброява типовете, които имат значение за
възстановяването.

\begin{table}[H]
\caption{Типове NAL единици в H.264, важни за възстановяването (по \cite{itu_h264})}
\label{tab:nal}
\centering
\begin{tabular}{@{}cL{4.2cm}L{7.0cm}@{}}
\toprule
\textbf{Тип} & \textbf{Наименование} & \textbf{Значение за възстановяването} \\
\midrule
1 & Част от кадър, който не е IDR & Декодира се само ако е наличен кадърът, от който зависи \\
5 & Част от IDR кадър & Декодира се самостоятелно, след него зависимостите се нулират \\
6 & SEI & Допълнителна информация, не е нужна за декодиране \\
7 & SPS & Параметри на поредицата: размер на кадъра, профил \\
8 & PPS & Параметри на картината, сочи към SPS \\
\bottomrule
\end{tabular}
\end{table}

Компресията в H.264 използва зависимости между кадрите \cite{wiegand2003}. Кадрите от
тип I се кодират само с информация от самите тях. Кадрите от тип P и B се кодират като
разлика спрямо други, вече декодирани кадри. Специален вид I кадър е \term{IDR кадърът}
(Instantaneous Decoding Refresh): след него нито един следващ кадър не се позовава на
кадрите преди него. Следствието за експертизата е пряко. Ако е загубен IDR кадър, всички
зависими от него кадри до следващия IDR кадър не могат да бъдат декодирани коректно.
Обратно, всеки оцелял IDR кадър е отправна точка, от която декодирането може да
започне наново.

Стандартът описва и два начина за разграничаване на NAL единиците в поток. В
\term{приложение B} на H.264 всяка единица се предшества от \term{стартов код}, байтовете
\code{00 00 01} \cite{itu_h264}. Във файл MP4 се използва друг запис: всяка единица се
предшества от полето ѝ за дължина, а стартови кодове няма \cite{iso14496_15}. Размерът
на това поле, както и SPS и PPS, се пазят в кутията \code{avcC} вътре в \code{stsd}, т.е.
отново в \code{moov}. Затова видеоданни от MP4 без \code{moov} не могат да се намерят
с просто търсене на стартови кодове. Методите в раздел~III трябва да разпознаят
полетата за дължина и да възстановят параметрите на кодека по друг път.
